\begin{textsample}{NEG Dim 4 – 2015\_c5 – Score 1.00 – 2015\_c5/enade\_2015\_tec\_mark\_6b65a1aace.txt}  \label{ex:f4_neg_185}
CAPÍTULO 5 DISTRIBUIÇÃO DOS CONCEITOS 5.1 PANORAMA NACIONAL DA DISTRIBUIÇÃO DOS CONCEITOS O Conceito Enade é um indicador de qualidade que avalia o desempenho dos estudantes a partir dos resultados do Enade.
Os valores possíveis vão de 1 ( pior situação ) a 5 ( melhor situação ).
As Unidades de Observação com apenas um ou sem nenhum concluinte participante não obtêm o Conceito Enade, ficando Sem Conceito ( SC ).
A Tabela 5.1 apresenta a quantidade e distribuição de cursos de Tecnologia em Marketing participantes do Enade/2015, por faixa de conceito e Grande Região.
Mais uma vez, ressaltando que a diferença entre os cursos tabulados neste capítulo e no capítulo 2 são os cursos sem conceito, em princípio, aqueles sem alunos concluintes que participassem da prova ou que tivessem tão somente um aluno concluinte¹⁸.
Observando -se os dados da Tabela 5.1, nota -se que, dos 271 cursos participantes, 106 ( 39,1 \% ) classificaram -se com conceito 3, o valor modal.
Este foi também o conceito modal em todas as regiões do país.
O conceito 2 foi o segundo mais frequente em nível nacional ( 21,0 \%, correspondendo a 57 cursos ) e o conceito 4, o terceiro ( 19,6 \%, correspondendo a 53 cursos ).
Houve, ainda, 13 cursos ( 4,8 \% ) que receberam conceito 5 e dez ( 3,7 \% ) cursos que receberam conceito 1.
Dos 271 cursos de Tecnologia em Marketing, 32 cursos ( 11,8 \% ) ficaram sem conceito ( SC ). ¹⁸ Estes cursos com somente um concluinte não foram considerados no capítulo 2.
A região Norte participou com sete cursos ou 2,6 \% do total nacional.
Nesta região, o conceito 3 foi atribuído a três cursos dos sete participantes, o que equivale a 42,9 \% do total regional, além de ser o conceito modal na região.
Nenhum curso ficou com conceito 4.
Três cursos foram avaliados com conceitos 1, 2 e 5 ( um curso em cada conceito ), correspondendo a 14,3 \%, cada.
Além disso, um curso ficou sem conceito ( SC ).
A região Nordeste participou com 37 cursos ou 13,7 \% do total nacional.
Nesta região, 13 cursos, 35,1 \% em termos regionais, obtiveram conceito 3, o conceito modal para a região.
Os conceitos 2 e 4 foram atribuídos a, respectivamente, dez e sete cursos ( 27,0 \% e 18,9 \% ).
Dois cursos ( 5,4 \% ) obtiveram o conceito 5 e um curso ( 2,7 \% ), o conceito 1.
Quatro cursos ( 10,8 \% ) ficaram sem conceito ( SC ) nesta região.
Dos 167 cursos participantes da região Sudeste ( 61,6 \% do total nacional ), o conceito 3 também foi o conceito modal, e foi obtido por 65 cursos ( 38,9 \% ).
O conceito 1 foi atribuído a oito cursos ( 4,8 \% ) e os conceitos 2 e 4, respectivamente, a 36 e 30 cursos, correspondendo a 21,6 \% e 18,0 \%.
O conceito 5 foi atribuído a sete cursos ( 4,2 \% ), e 21 cursos ( 12,6 \% ) ficaram sem conceito ( SC ).
Dos 50 cursos da região Sul, 20 ( 40 \% do total regional ) obtiveram o conceito 3 ( conceito modal ).
O conceito 2 foi atribuído a oito cursos ( 16,0 \% ), o conceito 4, a 16 cursos ( 32,0 \% ) e o conceito 5, a dois cursos ( 4,0 \% ).
Nenhum curso obteve conceito 1 e quatro cursos ( 8,0 \% ) na região Sul ficaram sem conceito ( SC ).
Já dos dez cursos participantes na região Centro-Oeste, cinco receberam conceito 3, o conceito modal ( 50,0 \% ).
Nesta região, dois cursos ficaram sem conceito ( SC ), enquanto nenhum ficou com conceitos 1 e 4.
Dois cursos foram \textbf{alocados} no conceito 2, correspondendo a 20,0 \% dos cursos nessa região, e um curso foi \textbf{alocado} ao conceito 5 ( 10,0 \% ).
O Gráfico 5.1 apresenta a distribuição do Conceito Enade segundo Grande Região : poligonais mais à esquerda representam Grandes Regiões cuja distribuição apresenta menores conceitos, e poligonais mais à direita, Grandes Regiões com melhores conceitos.
Apesar das poligonais entrecruzarem -se, pode -se observar que a região Norte ( linha azul ) apresenta os cursos em IES com conceitos tipicamente mais baixos, ao passo que a região Sul ( linha roxa ) apresenta a distribuição com valores maiores.
Cursos em IES das regiões Sudeste ( linha vermelha ), Centro-Oeste ( linha preta ) e Nordeste ( linha verde ) apresentam uma situação intermediária. 5.2 CONCEITOS POR CATEGORIA ADMINISTRATIVA E POR GRANDE REGIÃO A Tabela 5.2 apresenta a distribuição dos cursos participantes do Enade/2015 de Tecnologia em Marketing, por Categoria Administrativa, de acordo com os conceitos por eles alcançados, segundo as Grandes Regiões brasileiras.
Dos 271 cursos participantes, três ( 1,1 \% ) eram ministrados em Instituições Públicas e 268 ( 98,9 \% ), em Instituições Privadas.
De acordo com as informações da Tabela 5.2, em termos nacionais, todas as 13 IES que receberam o conceito 5 eram de Instituições Privadas.
Os dois cursos participantes de IES Públicas foram \textbf{alocados} no conceito 4.
Na rede privada, o conceito modal foi 3, com 106 cursos dos 268 da categoria.
Entre os demais cursos participantes, dez receberam conceito 1, e 57, conceito 2.
O Na análise por região, observa -se que as Instituições Públicas só tiveram participação nas regiões Sudeste e Sul.
Na região Norte, as Instituições Privadas participaram com todos os sete cursos desta região, dos quais um ficou sem conceito ( SC ) e três ficaram com conceitos 1, 2 e 5 ( um curso em cada conceito ).
O conceito modal para as Instituições Privadas na região foi 3, com três cursos, correspondendo a 42,9 \%.
Nenhum curso obteve conceito 4.
Na região Nordeste, a rede privada concentrou os 37 cursos participantes.
Dentre os cursos oferecidos por IES Privadas no Nordeste, 13 cursos ficaram \textbf{alocados} no conceito 3 ( conceito modal ), um curso obteve conceito 1 e dez, conceito 2.
Além disso, sete cursos receberam conceito 4 e dois, conceito 5.
Quatro cursos ficaram sem conceito ( SC ).
Na região Sudeste, 166 dos 167 cursos estavam concentrados na rede Privada, correspondendo a 99,4 \% do total regional.
Nesta categoria, na região Sudeste, o conceito modal foi 3 ( 65 cursos ).
Nesta combinação de Categoria Administrativa e Grande Região, 20 cursos não receberam conceito ( SC ).
Os demais foram avaliados com conceito 1 ( oito cursos ), conceito 2 ( 36 cursos ), conceito 4 ( 30 cursos ) e conceito 5 ( sete cursos ).
As Instituições Públicas participaram somente com um único curso, que ficou sem conceito ( SC ).
As Instituições Privadas concentraram 48 dos 50 cursos participantes da região Sul, 96,0 \% do total regional.
Desses, 20 ficaram com conceito 3, o conceito modal.
Nesta combinação de Categoria Administrativa e Grande Região, oito cursos obtiveram conceito 2, 14 cursos ficaram com conceito 4, enquanto dois cursos receberam conceito 5.
Nenhum curso foi associado ao conceito 1, enquanto quatro cursos ficaram sem conceito ( SC ).
As Instituições Públicas na região Sul participaram com dois cursos ( 4,0 \% do total regional ), e estes foram \textbf{alocados} ao conceito 4.
Na região Centro-Oeste, todos os dez cursos participantes também eram de Instituições Privadas.
Destes, cinco concentraram -se no conceito 3, conceito modal.
Dentre os demais, dois receberam conceito 2, e um recebeu conceito 5.
Dois cursos ficaram sem conceito ( SC ).
O Gráfico 5.2 apresenta a distribuição do Conceito Enade segundo a Categoria Administrativa da IES.
Os cursos em IES Públicas ( linha azul ) apresentam uma poligonal mais à direita do que os em IES Privadas ( linha verde ) e, consequentemente, uma distribuição de Conceitos Enade com valores maiores.
Cumpre notar que os cursos com notas em Instituições Públicas eram somente dois. 5.3 CONCEITOS POR ORGANIZAÇÃO ACADÊMICA E POR GRANDE REGIÃO Na Tabela 5.3, encontra -se a distribuição dos conceitos atribuídos aos cursos participantes do Enade/2015 na Área de Tecnologia em Marketing, por Organização Acadêmica, segundo as Grandes Regiões brasileiras.
Dos 271 cursos de Tecnologia em Marketing participantes, 81 eram oferecidos em Universidades, 56 cursos em Centros Universitários e 134 em Faculdades.
Esta distribuição corresponde a, respectivamente, 29,9 \%, 20,7 \% e 49,4 \% dos cursos.
De acordo com os dados apresentados, dos 13 cursos avaliados com conceito 5, cinco eram vinculados a Universidades.
Ficaram sem conceito ( SC ), 13 cursos oferecidos em Universidades.
Esse tipo de Organização Acadêmica teve o conceito 3 como modal, com 34 cursos ( 42,0 \% ).
Os demais cursos de Universidades avaliados receberam os conceitos 1 ( dois cursos ), 2 ( nove cursos ) e 4 ( 18 cursos ).
Entre os cursos em Centros Universitários, o conceito modal foi 3, com 25 cursos ( 44,6 \% ).
Neste tipo de Organização Acadêmica, apenas um curso ficou sem conceito ( SC ), e os outros cursos receberam os conceitos 1 ( um curso ), 2 ( 12 cursos ), 4 ( 14 cursos ) e 5 ( três cursos ).
Nas Faculdades, 18 dos 134 cursos ficaram sem conceito ( SC ), e 47 ( 35,1 \% ) ficaram \textbf{alocados} no conceito 3, conceito modal, seguido de 36 cursos que obtiveram conceito 2 e de 21 cursos que receberam conceito 4.
Os demais cursos neste tipo de Organização Acadêmica, receberam os conceitos 1 ( sete cursos ) e 5 ( cinco cursos ).
Considerando -se separadamente as regiões brasileiras, verifica -se que, na região Norte, as Universidades participaram com apenas um curso, e este foi \textbf{alocado} ao conceito 5.
Os Centros Universitários também participaram com um único curso, que foi \textbf{alocado} ao conceito 3.
As Faculdades concentraram cinco dos sete cursos na região Norte, dois receberam conceito 3 e dois obtiveram conceitos 1 e 2 ( um curso em cada conceito ).
Além disso, um curso ficou sem conceito ( SC ).
Na região Nordeste, as Universidades participaram com cinco dos 37 cursos na Área de Tecnologia em Marketing.
O conceito modal foi 3, com três cursos, enquanto os conceitos 4 e 5 foram atribuídos a um curso cada.
Nenhum dos cursos oferecidos por Universidades no Nordeste ficou sem conceito ( SC ) ou obteve os conceitos 1 ou 2.
Os Centros Universitários contaram com oito cursos participantes na região Nordeste, dos quais três obtiveram conceito 3 ( conceito modal ), quatro ficaram com conceitos 2 e 4 ( dois cursos em cada conceito ) e um curso foi \textbf{alocado} ao conceito 5.
Nenhum curso ficou sem conceito ( SC ) ou com conceito 1.
As Faculdades foram representadas por 24 cursos na região Nordeste, dos quais oito receberam conceito 2 ( conceito modal ), seguidos de sete cursos que receberam conceito 3 e quatro cursos com conceito 4.
Quatro cursos ficaram sem conceito ( SC ) e um curso foi \textbf{alocado} ao conceito 1.
Na região Sudeste, as Universidades concentraram 57 dos 167 cursos da região.
Entre os cursos em Universidades, nessa região, o conceito modal foi 3 com 24 cursos, enquanto 12 cursos ficaram sem conceito ( SC ).
Os demais cursos receberam os conceitos 1 ( dois cursos ), 2 ( sete cursos ), 4 ( 11 cursos ) e 5 ( um curso ).
Os Centros Universitários participaram com 39 cursos na região Sudeste, dos quais 17 obtiveram o conceito modal 3, enquanto um curso ficou sem conceito ( SC ).
Os demais receberam os conceitos 1 ( um curso ), 2 ( nove cursos),4 ( nove cursos ) e 5 ( dois cursos ).
As Faculdades foram representadas por 71 cursos na região Sudeste dos quais 24 receberam o conceito modal 3.
Os demais se distribuíram nos conceitos 1 ( cinco cursos ), 2 ( 20 cursos ), 4 ( dez cursos ) e 5 ( quatro cursos ).
Oito cursos ficaram sem conceito ( SC ) nessa categoria e região.
Dos 50 cursos da região Sul, 16 eram de Universidades, para os quais os conceitos modais foram 3 e 4 ( cada um desses conceitos atribuídos a seis cursos ).
Nesse tipo de organização, um dos cursos ficou sem conceito ( SC ) e nenhum, com o conceito 1.
Os demais receberam os conceitos 2 ( um curso ) e 5 ( dois cursos ).
Os Centros Universitários da região Sul participaram com oito cursos.
Destes, quatro cursos ficaram com conceito 3 ( conceito modal ), um ficou com conceito 2 e três com conceito 4.
Nenhum curso ficou sem conceito ( SC ) ou com os demais conceitos nesta organização acadêmica e região.
Foram 26 os cursos vinculados a Faculdades na região Sul, e dez desses receberam conceito 3, o conceito modal.
Seis cursos receberam o conceito 2 e sete o conceito 4.
Nenhum curso ficou com conceito 1 ou 5.
Três cursos ficaram sem conceito ( SC ).
Na região Centro-Oeste, dois dos dez cursos eram de Universidades, e estes foram \textbf{alocados} nos conceitos 2 e 3 ( um curso em cada ).
Os Centros Universitários da região Centro-Oeste não tiveram participação.
Dos oito cursos oferecidos em Faculdades na região Centro-Oeste, quatro ficaram \textbf{alocados} no conceito modal 3, e dois cursos obtiveram conceitos 2 e 5 ( um curso em cada ).
Dois cursos ficaram sem conceito ( SC ).
O Gráfico 5.3 apresenta a distribuição do Conceito Enade segundo a Organização Acadêmica da IES.
Os cursos em Universidades ( linha azul ) apresentam uma poligonal mais à direita do que as demais, denotando uma distribuição de Conceitos Enade melhores.
No outro extremo, temos os cursos oferecidos em Faculdades ( linha vermelha ) com a poligonal mais à esquerda e, consequentemente, uma distribuição de Conceitos Enade com menores valores.
IES de Centros Universitários ( linha verde ) apresentam uma situação intermediária.

% matched lemmas: alocar
\end{textsample}
